\chapter{Notation summary}\label{app:notation}
This appendix collects the symbols used across the thesis and, where one
exists, the equation in which each symbol is introduced. The symbols of the
piano setting come first, followed by those of the strings-and-winds setting
and of the waveform setting. Several letters are used in more than one role,
and the conventions that keep those roles apart are stated after the table.
\begin{center}
\begin{longtable}{@{}p{4.3cm}p{9.2cm}@{}}
\toprule
Symbol & Meaning \\
\midrule
\endfirsthead
\toprule
Symbol & Meaning \\
\midrule
\endhead
\bottomrule
\endfoot
$\score = \{s_i\}$, $s_i = (p_i, b_i, d_i)$ & score support: pitch, beat onset, beat duration \\
$y_i = [\tau_i, \log r_i, v_i]$ & piano variables: onset residual, articulation, dynamics \\
$T(\cdot)$ & tempo warp, score beats $\to$ performance seconds \\
$\graph$, $W$, $\Lg$ & score graph, adjacency \eqref{eq:adjacency}, Laplacian \eqref{eq:laplacian} \\
$\ell_b$, $\ell_p$, $w_v$ & graph length scales (beats, semitones), voice bonus \\
$\Qg$ & graph prior precision: $\lambda I + \eta \Lg$ or $\sigma_g^{-2}(\kappa^2 I + \Lg)^\alpha$ \\
$\lambda$, $\eta$ & additive ridge term, Laplacian weight \\
$\sigma_g$, $\kappa$, $\alpha$ & Mat\'ern scale, inverse length, exponent \\
$\mu$, $\muLM$, $h_i$ & prior mean; music-model mean; music-model per-note embedding \\
$B$, $c_f$ & coregionalization matrix; per-channel feature-kernel scales \eqref{eq:gpfirst} \\
$g(\nu; s)$, $s$ & shape-normalized spectral kernel of $\Lg$ ($g(0){=}1$); its shape parameter \\
$\nu$, $U$, $K_G$ & Laplacian eigenvalues and eigenvectors; graph kernel $U \operatorname{diag}(g(\nu;s)) U^\top$ \eqref{eq:spectral} \\
$d_5$, $\pi(p)$ & circle-of-fifths distance; pitch-class position map \eqref{eq:fifths}, used by \eqref{eq:tonaladj} \\
$K_{oo}$, $C$ & observed block of $K$ (noise included); latent covariance \eqref{eq:posterior} \\
$\beta$ & feature weights of the marginalized Bayesian linear mean (\S\ref{sec:features}) \\
$X_f$, $\varsigma^2$ & per-note feature matrices; per-channel observation noise \eqref{eq:gpfirst} \\
$\tilde y$, $e$, $\Sigma_e$ & observations, observation noise, its covariance \\
$m$, $\Sy$, $\sigma_i$ & posterior mean, posterior covariance, posterior standard deviation ($\sigma$ has no other use in this thesis) \\
$y^{(\zeta)}$, $C^{(\zeta)}$, $m^{(\zeta)}$ & prior components ($\zeta \in \{G, \mathrm{feat}, \mathrm{emb}\}$ on piano), their covariances, posterior-mean components \eqref{eq:components}--\eqref{eq:decomp} \\
$\omega^{(\zeta)}$, $\rho^{(\zeta\zeta')}_i$ & covariance share \eqref{eq:share}; per-note component correlation \eqref{eq:compcorr} \\
$\sigma_i^{\mathrm{pred}}$ & floored predictive std, $\sqrt{(\Sy)_{ii} + \sigma_e^2}$ \\
$c_i$, $f_0$ & intonation (cents), fundamental frequency (strings and winds) \\
$\ell_i$, $f_i^{\mathrm{vib}}$, $\gamma_i$, $d_i^{\mathrm{vib}}$ & strings-and-winds loudness (log RMS); vibrato rate, extent, onset delay \\
$\sigma_{e,ij}^2$ & per-(note, channel) observation-noise variance (strings and winds) \\
$J_i$, $\pi_j$, $q_{0.2}$ & note-$i$ frame set; tracker per-frame confidence; its lowest-quintile threshold \eqref{eq:framerule} \\
$\hat\varphi_i$, $\hat\psi_i$ & local tempo-line slope and intercept at note $i$ \eqref{eq:localwarp} \\
$f_{\mathrm{ref}}$ & tuning reference, $440$\,Hz $=$ A4 (MIDI $69$) \eqref{eq:cents-curve} \\
$x$, $z$, $a$, $\Phi(z)$, $\epsilon$, $K_n$, $H$ & audio, nonlinear positions, amplitudes, synth basis, noise, noise cov, harmonic count (the waveform) \\
$A_i(t)$ & amplitude envelope of note $i$ (the waveform) \\
\end{longtable}
\end{center}
\noindent A first group of symbols is kept apart by a fixed rule. The
support $\score$ is never a covariance. Plain $\sigma$ is never a prior
scale, whereas $\sigma_g$ is one. The exponent $\alpha$ is never the
additive weight, which is $\eta$. The capital $A_i(t)$ is the amplitude
envelope, and the lowercase $a_i(t)$ is never used for it. $\Phi(z)$ is the
synthesis operator, while $\Phi_{\N}$ is the normal cumulative distribution
function.

A second group of overloads is resolved by context. The letter $s$ is the
spectral shape parameter, while $s_i$ is a score event. Capital $P$ is the
per-note observation precision. The map $\pi(p)$ is the circle-of-fifths
position, while $\pi_j$ is a tracker confidence. $B$ is the
coregionalization matrix, and $B = 2000$ appears only as the bootstrap
replicate count. $\delta_i$ is the performed duration and
$\delta_{\zeta\zeta'}$ the Kronecker delta, and the vibrato onset delay is
written $d_i^{\mathrm{vib}}$. The lengths $\ell_b, \ell_p$ are graph length
scales, while $\ell_i$ is the loudness channel. The letter $c$ is a channel
index, $c_f$ a feature-kernel scale, and $c_i$ the intonation channel in
cents. An entry of the strings-and-winds mask is a (note, channel) pair,
whereas the evaluation unit of the piano tables is a (piece, seed)
combination and is always written out as a fit. Finally, the exponent
$\alpha$ and the eigenvalue $\nu$ of this thesis differ from the Mat\'ern
smoothness convention of the graph-GP literature.

% =====================================================================
